%%%%%%%%%%%%%%%%%%%%%%%%%%%%%%%%%%%%%%%%%%%%%%%%%%%%%%
%% "Magic comments" for the LaTeX editor -------------
% !TEX root = ../Diss.tex
%%%%%%%%%%%%%%%%%%%%%%%%%%%%%%%%%%%%%%%%%%%%%%%%%%%%%%
%
%%%%%%%%%%%%%%%%%%%%%%%%%%%%%%%%%%%%%
\addchap{Kurzeinführung in diese Vorlage}
%%%%%%%%%%%%%%%%%%%%%%%%%%%%%%%%%%%%%
%
%% Add an entry for this chapter to the table of contents:
%\addcontentsline{toc}{chapter}{Kurzeinführung in die Vorlage}
%
%% If no entry to the table of contents:
%% Set a PDF bookmark (PDF-Lesezeichen):
%\pdfbookmark[0]{Kurzeinführung in die Vorlage}{kurzeinfuehrung}
%
%Add a header title (Kolumnentitel beidseitig setzen)
\markboth{Kurzeinführung in die Vorlage}{Kurzeinführung in die Vorlage}
%
%%%%%%%%%%%%%%%%%%%%%%%%%%%%%%%%%%%%%
%
Herzlich Willkommen zur \LaTeX-Vorlage
%Nutzung von \glsdisp wegen der Genitiv-Form
des \glsgen{ac:IES} %{Lehrstuhls für Interalktive Echtzeitsysteme (IES)}
des \glsgen{ac:KIT}. %{Karlsruher Institus für Technilogie (KIT)}.
Sie wurde ursprünglich von
\href{https://ies.iar.kit.edu/1473_1575.php}{Philipp Woock}
entwickelt und basiert in ihren Grundzügen auf der \enquote{allgemeinen sehr umfassenden}
\href{http://www.matthiaspospiech.de/latex/templates/thesis/}{Vorlage} von 
\href{http://www.matthiaspospiech.de}{Matthias Pospiech} von der Leibniz Universität Hannover.
Ohne diese Basis wäre die Vorlage niemals das geworden, was sie ist. Vielen Dank!

Im Jahre 2018 wurde diese Vorlage von
\href{https://ies.iar.kit.edu/1473_1650.php}{Michael Grinberg}
grundlegend überarbeitet und an die neuen Vorgaben des \glsgen{ac:KSP} angepasst.
Insbesondere wurde die Vorlage dahingehend abgeändert, dass sie nun XeLaTeX statt pdfLaTeX, 
\gls{gls:biblatex} statt \gls{gls:bibtex} verwendet und die inzwischen aufgetretenen
Inkompatibilitäten mit neuen Versionen einiger \gls{gls:latex}-\glspl{gls:pkg} behebt.
In 2019 wurde die Vorlage weiter angepasst und mit Beispielen aus einer Vorlage
für studentische Arbeiten, die von
\href{https://crypto.iti.kit.edu/alumni_matthias_nagel.php}{Mathias Nagel}
entwickelt worden ist, angereichert.

Insgesamt erfüllt diese Vorlage nun einen doppelten Zweck:
zum einen gibt sie eine standardisierte und mit dem \gls{ac:KSP}
abgestimmte Formatierung vor, zum anderen demonstriert sie diverse
standardmäßige und eigene Funktionalitäten und enthält zudem
eine ausführliche Anleitung zur Nutzung dieser.

Die eigentliche Anleitung befindet sich in \cref{chap:Examples} (\nameref{chap:Examples}).
Texte in den Kapiteln \nameref{chap:abstract}, \nameref{chap:acknowledgements} und \nameref{chap:kurzfassung} sind Platzhalter.

Zur produktiven Nutzung dieser Vorlage empfiehlt es sich, 
die Anleitungskapiteln auszublenden.
Dies geschieht in der Datei
\printfilepathquoted{preambel/AlleSchalter.tex} in der Zeile
\lstinline|\setUserDefinedBoolean{showExamples}{true}|. Hier soll 
der Wert \printkeyword{true} durch \printkeyword{false} ersetzt werden.
Zur Einblendung einer vordefinierten Kapitelstruktur einer Dissertation
sollte stattdessen der Wert der Variablen
\printkeyword{showIntro},
\printkeyword{showRelatedWork},
\printkeyword{showConcept}
\usw
%\printkeyword{showSystem},
%\printkeyword{showEvaluation},
%\printkeyword{showConclusion} und
%\printkeyword{showOutlook}
auf \printkeyword{true} gesetzt werden.



\begin{tcolorbox}[colback=orange]
\footnotesize
Hinweis:
Zu beachten ist, dass zum Kompilieren der Vorlage
\printkeyword{xelatex} statt \printkeyword{pdflatex}
und zum Kompilieren der Bibliographien
\printkeyword{biber} statt \printkeyword{bibtex}
verwendet werden sollen.
Zur Generierung des Begriffs- und Abkürzungsverzeichnisses muss zusätzlich
\printkeyword{makeglossaries} (benötigt Perl!) aufgerufen werden.
\end{tcolorbox}
